\abstract{
\thispagestyle{fancy}

\begin{otherlanguage}{ukrainian}

\textit{Корнієнко С.С.} Методика навчання цифрової історії майбутніх фахівців. -- Кваліфікаційна наукова праця на правах рукопису.

Дисертація на здобуття наукового ступеня доктора філософії за спеціальністю 011 Освітні, педагогічні науки. -- Криворізький державний педагогічний університет, Кривий Ріг, 2026.

{Дисертаційну роботу присвячено теоретичному обґрунтуванню, розробці та експериментальній перевірці методики навчання цифрової історії майбутніх фахівців. Дослідження спрямоване на вирішення наукової проблеми розробки методики навчання цифрової історії в закладах вищої освіти, яка забезпечує формування цифрових компетентностей та навичок критичного мислення.

У першому розділі подано теоретичні засади дослідження. За результатами систематичного огляду за методологією PRISMA було ідентифіковано 6 педагогічних досліджень, присвячених цифровій історії в освіті, що свідчить про початковий стан розробленості цієї проблеми. Систематичний огляд 74 досліджень гейміфікації в навчанні історії виявив, що лише 14,9\% із них пропонують чітке визначення гейміфікації, 83,8\% мають високий ризик упередженості за методологією MRQS, а 85,1\% не надають операціональних визначень ключових конструктів. Аналіз якості досліджень за допомогою MRQS з валідацією великими мовними моделями (LLM) підтвердив ці результати та продемонстрував потенціал LLM для автоматизації оцінки якості досліджень. Розроблено рекомендації щодо впровадження гейміфікації на основі рамки цифрових компетентностей DigComp~2.0.

У другому розділі представлено методологію педагогічного експерименту. Дослідження проводилось на базі Центру довузівської підготовки та освіти Криворізького державного педагогічного університету протягом жовтня 2024~-- квітня 2025~року. У курсі підготовки до НМТ з історії України взяли участь 23 слухачі. Навчальний процес включав 38 лекцій, 30 семінарів із 156 тестовими формами Google Forms, 5 модульних контрольних робіт та інтерактивну гейміфіковану платформу HistoryLikeGame з 23 ігровими модулями. Збір даних здійснювався за допомогою Google Forms, журналу чату Viber та аналітики гейміфікованої платформи.

У третьому розділі подано результати педагогічного експерименту та їх обговорення. Аналіз відповідей студентів на семінарах 25--40 продемонстрував діапазон середніх результатів від 33,3\% до 83,3\%. Відслідковано індивідуальні траєкторії навчання 16 учасників. Аналіз відвідуваності виявив 33 повідомлення про відсутність, основними причинами яких стали проблеми з інтернетом (24\%), сімейні обставини (24\%) та відключення електроенергії (9\%)~-- типові для умов воєнного часу. Гейміфікована платформа, впроваджена 1 грудня 2024 року, доповнила традиційні методи оцінювання інтерактивними елементами drag-and-drop та текстового введення.

Наукова новизна отриманих результатів полягає в тому, що:
\textit{вперше} розроблено та теоретично обґрунтовано методику навчання цифрової історії майбутніх фахівців, яка інтегрує гейміфікацію як ключовий компонент цифрової педагогіки; проведено систематичний огляд досліджень гейміфікації в навчанні історії за методологією PRISMA (74 дослідження) та оцінку їхньої якості за MRQS; здійснено валідацію автоматизованої оцінки якості досліджень за допомогою великих мовних моделей;
\textit{удосконалено} зміст та методи навчання історії з використанням цифрових засобів у закладах вищої освіти;
\textit{набули подальшого розвитку} теоретичні засади застосування гейміфікації та цифрових технологій у навчанні гуманітарних дисциплін.

Практичне значення отриманих результатів полягає в тому, що розроблено інтерактивну гейміфіковану платформу HistoryLikeGame та комплект навчально-методичних матеріалів для підготовчих курсів з історії України.

\textit{Ключові слова}: цифрова історія, гейміфікація, методика навчання історії, систематичний огляд, PRISMA, MRQS, цифрова компетентність, DigComp~2.0, підготовка до НМТ, дистанційне навчання, освіта в умовах воєнного часу.}

\end{otherlanguage}

\begin{center}
\large
    \textbf{ABSTRACT}
\end{center}

\textit{Korniienko S.S.} Digital history teaching methodology for future professionals. -- Qualifying scientific work on manuscript rights.

Dissertation for the Doctor of Philosophy degree in the speciality 011 Education and Pedagogical Sciences. -- Kryvyi Rih State Pedagogical University, Kryvyi Rih, 2026.

{This dissertation is devoted to the theoretical substantiation, development, and experimental verification of a digital history teaching methodology for future professionals. The research addresses the scientific problem of developing a methodology for teaching digital history in higher education institutions that ensures the formation of digital competencies and critical thinking skills.

Chapter~1 presents the theoretical foundations. A PRISMA-guided systematic review identified only 6 pedagogical studies dedicated to digital history in education, indicating the nascent state of the field. A systematic review of 74 studies on gamification in history education revealed that only 14.9\% provided a clear definition of gamification, 83.8\% carried a high risk of bias according to the MRQS methodology, and 85.1\% failed to provide operational definitions of key constructs. A quality assessment using MRQS with large language model (LLM) validation confirmed these findings and demonstrated the potential of LLMs for automating research quality evaluation. Guidelines for implementing gamification based on the DigComp~2.0 digital competence framework were developed.

Chapter~2 describes the experimental methodology. The study was conducted at the Centre for Pre-University Training and Education of Kryvyi Rih State Pedagogical University during October 2024 -- April 2025. Twenty-three students participated in an NMT (National Multi-Subject Test) preparatory course in Ukrainian history. The educational process comprised 38 lectures, 30 seminars with 156 Google Forms quizzes, 5 modular controls, and an interactive gamified platform HistoryLikeGame with 23 game modules. Data were collected via Google Forms, a Viber chat log, and gamification platform analytics.

Chapter~3 presents the results and their discussion. Analysis of student responses across Seminars 25--40 demonstrated mean scores ranging from 33.3\% to 83.3\%. Individual learning trajectories were tracked for 16 participants. Attendance analysis revealed 33 absence reports, with primary reasons including internet issues (24\%), family circumstances (24\%), and power outages (9\%) -- characteristic of wartime conditions. The gamified platform, introduced on December 1, 2024, supplemented traditional assessment methods with interactive drag-and-drop and text-input elements.

The scientific novelty of the obtained results lies in the fact that:
\textit{for the first time,} a methodology for teaching digital history to future professionals that integrates gamification as a key component of digital pedagogy has been developed and theoretically substantiated; a PRISMA-guided systematic review of gamification studies in history education (74 studies) and their quality assessment using MRQS have been conducted; validation of automated research quality assessment using large language models has been performed;
\textit{improved:} the content and methods of history education using digital tools in higher education institutions;
\textit{further developed:} the theoretical foundations for applying gamification and digital technologies in humanities education.

The practical significance of the obtained results lies in the development of the interactive gamified platform HistoryLikeGame and a set of educational and methodological materials for Ukrainian history preparatory courses.

\textit{Keywords}: digital history, gamification, history teaching methodology, systematic review, PRISMA, MRQS, digital competence, DigComp~2.0, NMT preparation, distance learning, wartime education.}


\begin{center}
    \textbf{List of education seeker publications}
\end{center}


\printpratsi

} % end abstract
